\documentclass[10pt,a4paper]{article}
\usepackage[utf8]{inputenc}
\usepackage[a4paper,left=3cm,right=2cm]{geometry}
\usepackage{amsmath}
\usepackage[thmmarks, amsmath, thref]{ntheorem}
\usepackage{amsfonts}
\usepackage{amssymb}
\usepackage{fancyhdr}
\pagestyle{fancy}
\usepackage{anysize}
\usepackage{graphicx}
\usepackage{float}
\usepackage{hyperref}
\usepackage{enumitem}
\newtheorem{theorem}{Theorem}
\theoremsymbol{\ensuremath{\blacksquare}}
\newtheorem*{solution}{Solución:}

\usepackage{xcolor}

% Margenes y formalidades

\textwidth 6.25in % Margen Derecho y ancho del texto
\textheight 8.5in
\oddsidemargin 0in % Margen Izquierdo
\headheight 0in % Largo del texto
\leftmargin 1in
\parindent 0em % Salto entre la indentación (Sangría)https://www.overleaf.com/4519127669sdpptkhvsbfn
\topmargin -0.5in % Margen Superior
\parskip 0.5ex % Salto entre parrafos
\baselineskip 1pt



\lhead{
\includegraphics[scale=0.1]{figuras/logoUSM.png}} % texto izquierda de la cabecera
%\chead{}                                      % texto centro de la cabecera
\rhead{Universidad Técnica Federico Santa María\\Ingeniería Comercial } % número de página a la derecha
\renewcommand{\headrulewidth}{0.4pt} % grosor de la línea de la cabecera
\renewcommand{\footrulewidth}{0.4pt} % grosor de la línea del pie

\begin{document}
\vspace{5cm}

\begin{center}
    {\Large  Ejercicios 5: Econometría }

Profesor:    Pedro Comte Hidalgo
\end{center}



\section{Ejercicios}
\begin{enumerate}
    \item Supongamos un modelo de Regresión Lineal Simple, que describe el salario en función de la educación de la forma:
    $$salary = \beta_0 + \beta_1 educ$$
    Si la solución mediante MCO, los entrega la siguiente forma:
        $$\hat{salary} = 10 + 2\cdot educ + \hat{u}$$
    siendo $\hat{u}$ el error observado.
    \begin{enumerate}
        \item  ¿Como se interpreta el valor de $\hat{\beta_1}$?  \\
        \textbf{el valor nos indica que el salario aumenta en 2 por cada año adicional de educación}

        \item Si dada la gran escala de $salary$, se le aplica $log()$, para tener valores más faciles de trabajar, encontrando la siguente RLS mediante MCO:
        $$\hat{log({salary})} = 2.8 + 0.07\cdot educ + \hat{u}$$
        ¿Como se interpreta el valor de $\hat{\beta_1}$ en este caso?  \\
        \textbf{El valor nos indica que el salario aumenta 7\% por cada 1 año adicional de educación}
        \item Ahora, se le aplica $log()$ también a educ, para tener valores más faciles de trabajar, encontrando la siguente RLS mediante MCO:
        $$\hat{log({salary})} = 1.4 + 0.80\cdot log(educ) + \hat{u}$$
        ¿Como se interpreta el valor de $\hat{\beta_1}$ en este caso? ¿En economia, como se le llama a este valor?\\
        \textbf{$\hat{\beta_1}$ en esta caso nos indica el porcentaje de aumento del salario al aumentar 1\% la educación, es decir, aumenta 0.8\% al aumentar 1\% la educación. Este valor es conocido como la elasticidad.}
    \end{enumerate}


\item Se tiene un modelo que explica el puntaje en la PAES dado el IQ del individuo, de la sigueinte forma:
$$Puntaje = 200 + 10 \cdot IQ -  0.05IQ^2$$
\begin{enumerate}
    \item ¿Cual es el efecto marginal del IQ en el puntaje?
    \textbf{$$\frac{\Delta puntaje}{\Delta IQ} = 10-0.1IQ$$ }
    \item ¿Cual es el valor donde el efecto marginal del IQ es cero?
    \textbf{$10-0.1IQ=0 \implies IQ = 100$}
    \item ¿Que sucede para valores de IQ mayores a 100?
    \textbf{El efecto del IQ en el puntaje empieza a ser negativo}
\end{enumerate}

\item Se tiene la siguiente RLS:
$$y = \beta_0 + \beta_1x$$
mediante MCO podemos encontrar $\hat{\beta_0}$ y $\hat{\beta_1}$ insesgados y eficientes. Si reescalamos $\Tilde{x} = 10x$. 
\begin{enumerate}
    \item Demuestre que $\hat{\beta_1}=10\hat{\Tilde{\beta_1}}$. \\
    \textbf{La RLS podemos definirla como:
    $$\hat{y} = \hat{\beta_0} + \hat{\Tilde{\beta_1}}\Tilde{x}+u = \hat{y} = \hat{\beta_0} + \hat{\Tilde{\beta_1}}(10x)+u$$
    por MCO sabemos que:
    $$\hat{\Tilde{\beta_1}}=\frac{Cov(10x,y)}{Var(10x)}=\frac{10Cov(x,y)}{100Var(x)}=\frac{1}{10}\frac{Cov(x,y)}{Var(x)}$$
    y sabemos que:
    $$\frac{Cov(x,y)}{Var(x)} = \hat{\beta_1}$$
    entonces:
    $$\hat{\Tilde{\beta_1}} = \frac{1}{10}\hat{\beta_1} \implies \hat{\beta_1} = 10\hat{\Tilde{\beta_1}}$$}
\item Demuestra que el valor $t$ del estimador no cambia. \\
\textbf{ Sea el t del estimador:
$$t_{\hat{\beta_1}} = \frac{\hat{\beta_1}}{s.e.(\hat{\beta_1})}$$
y para $\hat{\Tilde{\beta_1}}$ se tiene:
$$t_{\hat{\Tilde{\beta_1}}} = \frac{\hat{\Tilde{\beta_1}}}{s.e.(\hat{\Tilde{\beta_1}})}$$
Notames que 
$$s.e.(\hat{\Tilde{\beta_1}}) = \sqrt{Var(\hat{\Tilde{\beta_1}})}=\sqrt{Var(\frac{1}{10}\hat{{\beta_1}})}=\frac{1}{10}\sqrt{Var(\hat{{\beta_1}})}$$
entonces:
$$t_{\hat{\Tilde{\beta_1}}} = \frac{\frac{1}{10}\hat{{\beta_1}}}{\frac{1}{10}\sqrt{Var(\hat{{\beta_1}})}} = t_{\hat{\beta_1}}$$
}
    
\end{enumerate}

\item Se tiene los siguentes resultados de un modelo de Regresión Lineal, de las notas finales de un curso, dada caracteristicas del alumno. De los 674 datos se tienen los siguientes resultados de los estimadores:

\begin{center}
    \includegraphics[scale=1]{Ejercicios/sem1-2/eje5.jpg}
\end{center}

\begin{enumerate}
\item Con un nivel de significancia del 5\%, defina cuales son las variables que son significativamente distintas de cero.
\\
\textbf{Primero se tiene que entrar el valor crítico de t (dos colas) para el nivel de significancia y los grados de libertad. Tenemos: 
$$t_{1-\alpha/2,n-k-1}=t_{0.975,668} = 1.96$$
Luego se encuentra el estadistico t para cada variable:
$$t_{attend} = \frac{-0.06926}{0.04074}=-1.7$$
$$t_{termpGPA}=\frac{3.54031}{0.31395}=11.28$$
$$t_{priGPA}=\frac{-0.06566}{0.38589} = -0.17$$
$$t_{ACT} = \frac{0.27214}{0.05038} = 5,40$$
$$t_{hwrte} = \frac{-0.01451}{0.01038} = -1,40$$
Comparando sus valores absolutos con el valor t crítico, tenemos que los estimadores de $termGPA$ y $ACT$, son significativamente distintos de cero, con un 5\% de significancia. Mientas que para $attend$, $priGPA$ y $hwrte$, no se rechaza la hipotesis nula, y no serían relevante para este modelo.
}

\item Haga el intervalo de confianza para el estimador de la variable $attend$, al 95\% de confianza. Concluya.\\
\textbf{ el IC esta dado por:
$$[\hat{\beta} - t_{1-\alpha/2,n-k-1}s.e(\hat{\beta});\hat{\beta} + t_{1-\alpha/2,n-k-1}s.e(\hat{\beta})]$$
por tanto reemplazando tenemos:
$$[-0.06926 - 1.96\cdot 0.04074;-0.06926 - 1.96\cdot 0.04074]$$
$$[-0.149;0.011]$$
Como ven el intervalo contiene el cero, por lo que con un 95\% el estimador podría ser cero.}

\item Haga el intervalo de confianza para el estimador de la variable $termGPA$, al 95\% de confianza. Concluya.\\
\textbf{ el IC esta dado por:
$$[\hat{\beta} - t_{1-\alpha/2,n-k-1}s.e(\hat{\beta});\hat{\beta} + t_{1-\alpha/2,n-k-1}s.e(\hat{\beta})]$$
por tanto reemplazando tenemos:
$$[3.54031 - 1.96\cdot 0.31395;3.54031 + 1.96\cdot 0.31395]$$
$$[2.925;4.156]$$
Con un 95\% de confianza el estimador de $termGPA$ está entre 2.925 y 4.156.}
\end{enumerate}

\item Se tiene los siguientes resultados de un modelo de 5 variables:
\begin{center}
    \includegraphics[scale=0.7]{Ejercicios/sem1-2/eje5.2.jpg}
\end{center}
Luego se sacan dos variables y se obtienen los siguientes resultados:
\begin{center}
    \includegraphics[scale=0.7]{Ejercicios/sem1-2/eje5.3.jpg}
\end{center}
Se quiere saber si las variables excluidas del modelo son relevantes para el mismo, con una significancia del 5\% con valor critico de $F_{95\%,2,668}=3.01$, siendo 2 las variables que se eliminan, y n-k-1=668

\textbf{
$$F_{estadistico} = \frac{R^2_{ur}-R^2_{r}}{1-R^2_{ur}}\frac{n-k-1}{j}$$
$$F_{estadistico} = \frac{0.333-0.1753}{1.0,333}\frac{668}{2}=36,82$$
Como $F_{estadistico}>F_{95\%,2,668}$ las variables excluidas si son relevantes.\\

}


\end{enumerate}

    
\end{document}